  \footnote{It is known that widely used software libraries have differing definitions of the queryname sort order, meaning care should be taken when operating on multiple files of varying provenance.
  Tools may wish to use the sub-sort field to explicitly distinguish between natural and lexicographical ordering.  See Section~\ref{sec:sub-sort}.}\\\cline{2-3}
  & {\tt GO} & Grouping of alignments, indicating that similar alignment records
    are grouped together but the file is not necessarily sorted overall.
    \emph{Valid values}: {\tt none} (default), {\tt query} (alignments are
    grouped by {\sf QNAME}), and {\tt reference} (alignments are grouped by
    {\sf RNAME}/{\sf POS}).\\\cline{2-3}
  & {\tt SS} & Sub-sorting order of alignments. Valid values are of the form \emph{sort-order}{\tt :}\emph{sub-sort}, where \emph{sort-order} is the same value stored in the {\tt SO} tag and \emph{sub-sort} is an implementation-dependent colon-separated string further describing the sort order, but with some predefined terms defined in Section~\ref{sec:sub-sort}.
    For example, if an algorithm relies on a {\tt coordinate} sort that, at each coordinate, is further sorted by query name then the header could contain {\tt @HD SO:coordinate SS:coordinate:queryname}.%
\footnote{The repetition of \emph{sort-order} enables a limited form of validation.
For example, {\tt @HD SO:queryname SS:coordinate:TLEN} would indicate that the data has been re-sorted (by query name) by a non-{\tt SS}-aware tool and the {\tt SS} field should be ignored.}
    If the primary sort is not one of the predefined primary sort orders, then {\tt unsorted} should be used and the sub-sort is effectively the major sort.
    For example, if sorted by an auxiliary tag {\tt MI} then by coordinate then the header could contain {\tt @HD SO:unsorted SS:unsorted:MI:coordinate}.
    \newline
    \emph{Regular expression}: {\tt (coordinate|queryname|unsorted)(:[\verb"A-Za-z0-9_-"]+)+}\\\cline{1-3}
  \multicolumn{2}{|l}{\tt @SQ} & Reference sequence dictionary. The order of {\tt @SQ} lines defines the alignment sorting order.\\\cline{2-3}
  & {\tt SN}* & Reference sequence name.
The {\tt SN} tags and all individual {\tt AN} names in all {\tt @SQ} lines
must be distinct.
  The value of this field is used in the
  alignment records in {\sf RNAME} and {\sf RNEXT} fields. Regular expression: {\tt\rnameRegexp}\\\cline{2-3}
  & {\tt LN}* & Reference sequence length. \emph{Range}: $[1,\,2^{31}-1]$\\\cline{2-3}
  & {\tt AH} & Indicates that this sequence is an alternate locus.%
\footnote{See \url{https://www.ncbi.nlm.nih.gov/grc/help/definitions} for descriptions of \emph{alternate locus} and \emph{primary assembly}.}
  The value is the locus in the primary assembly for which this sequence is an alternative, in the format `\emph{chr}{\tt :}\emph{start}{\tt -}\emph{end}', `\emph{chr}' (if known), or `{\tt *}' (if unknown), where `\emph{chr}' is a sequence in the primary assembly.
  Must not be present on sequences in the primary assembly.\\\cline{2-3}
  & {\tt AN} & Alternative reference sequence names.
A comma-separated list of alternative names that tools may use when referring
to this reference sequence.%
\footnote{For example, given `{\tt @SQ SN:MT AN:chrMT,M,chrM LN:16569 TP:circular}',
tools can ensure that a user's request for any of `MT', `chrMT', `M',
or~`chrM' succeeds and refers to the same sequence.}
These alternative names are not used elsewhere within the SAM file;
in particular, they must not appear in alignment records' {\sf RNAME}
or~{\sf RNEXT} fields.
\emph{Regular expression}: \emph{name}{\tt (,}\emph{name}{\tt )*}
where \emph{name} is {\tt\rnameRegexp}\\\cline{2-3}
  & {\tt AS} & Genome assembly identifier. \\\cline{2-3}
  & {\tt DS} & Description.  UTF-8 encoding may be used.\\\cline{2-3}
  & {\tt M5} & MD5 checksum of the sequence.  See Section~\ref{sec:ref-md5}\\\cline{2-3}
  & {\tt SP} & Species.\\\cline{2-3}
  & {\tt TP} & Molecule topology. \emph{Valid values}: {\tt linear} (default) and {\tt circular}.%
\footnote{The previous footnote's example identifies MT as a circular chromosome.
The {\tt TP} field is often omitted, which implies linear.}\\\cline{2-3}
  & {\tt UR} & URI of the sequence.  This value may start with one of the standard
  protocols, e.g., `{\tt http:}' or `{\tt ftp:}'.  If it does not start with one of these protocols, it is assumed to be a file-system path.\\\cline{1-3}
  \multicolumn{2}{|l}{\tt @RG} & Read group. Unordered multiple {\tt @RG} lines are allowed.\\\cline{2-3}
  & {\tt ID}* & Read group identifier. Each {\tt @RG} line must have a unique {\tt ID}. The value of {\tt ID}
  is used in the RG tags of alignment records. Must be unique among all read groups in header section.  Read group IDs may be modified when merging SAM files in order to handle collisions.\\\cline{2-3}
  & {\tt BC} & Barcode sequence identifying the sample or library.
  This value is the expected barcode bases as read by the sequencing machine in the absence of errors.
\end{description}

\subsubsection{Reference MD5 calculation}\label{sec:ref-md5}
The {\tt M5} tag on {\tt @SQ} lines allows reference sequences to be uniquely identified through the MD5 digest of the sequence itself.
As the digest is based on the sequence and nothing else, it can help resolve ambiguities with reference naming.
For example, it allows a quick way of checking that references named `1', `Chr1' and `chr1' in different files are in fact the same.

The reference sequence must be in the 7-bit US-ASCII character set.
All valid reference bases can be represented in this set, and it avoids the problem of determining exactly which 8-bit representation may have been used.
Padding characters (See Section~\ref{sec:padded-sam}) must be represented only using the `*' character. 

The digest is calculated as follows:
\begin{itemize}
\item All characters outside of the inclusive range 33 (`\char33') to 126 (`\char126') are stripped out.
This removes all unprintable and whitespace characters including spaces and new lines.
Everything else is retained, even if not a legal nucleotide code.
\item All lowercase characters are converted to uppercase.
This operation is equivalent to calling \href{http://pubs.opengroup.org/onlinepubs/9699919799/functions/toupper.html}{\sl toupper()} on characters in the \href{http://pubs.opengroup.org/onlinepubs/9699919799/basedefs/V1\_chap07.html\#tag\_07\_02}{\sl POSIX locale}.
\item The MD5 digest is calculated as described in \href{https://tools.ietf.org/html/rfc1321}{\sl RFC 1321} and presented as a 32 character lowercase hexadecimal number.
\end{itemize}

As an example, if the reference contains the following characters (including spaces):
